% ASSERT_CONVENTION: entropy_base=nats, generator_convention=probabilist_dpdt_pQ, matrix_norm=sup_norm_rows, experiential_density=I*(1-I/H)
%
% Theorem A: Complete Proof by Composition of Lemmas
% Phase: 01-theorem-a-assembly, Plan: 02 (revised)
%
% This document assembles the complete proof of Theorem A from the seven
% lemma statements established in Plan 01 (derivations/theorem-a-lemmas.tex).
% Every step references a specific lemma or performs explicit algebra.
%
% Revision history:
%   v2: Fixed mu_stable lower bound (case analysis for tau_exit vs T_eps),
%       corrected gamma = alpha/2 (honest value from eta choice),
%       added routing factor discussion for general p.
%
% References:
%   [BEGK02]  Bovier, Eckhoff, Gayrard, Klein, Comm. Math. Phys. 228 (2002) 219--255
%   [BdH15]   Bovier, den Hollander, Metastability: A Potential-Theoretic Approach, Springer (2015)
%   [FW12]    Freidlin, Wentzell, Random Perturbations of Dynamical Systems, 3rd ed., Springer (2012)
%   [CV16]    Champagnat, Villemonais, Probab. Theory Relat. Fields 164 (2016) 243--283
%   [DV75-83] Donsker, Varadhan, Comm. Pure Appl. Math. (1975--1983)

\documentclass[11pt]{article}
\usepackage{amsmath,amssymb,amsthm}
\usepackage[margin=1in]{geometry}

\newtheorem{theorem}{Theorem}
\newtheorem{lemma}{Lemma}
\newtheorem{remark}{Remark}
\newtheorem{proposition}{Proposition}

\newcommand{\eps}{\varepsilon}
\newcommand{\Ds}{\Delta_s}
\newcommand{\Db}{\Delta_b}
\newcommand{\Bstable}{B_{\mathrm{stable}}}
\newcommand{\BBB}{B_{\mathrm{BB}}}
\newcommand{\Teps}{T_\eps}
\newcommand{\TV}{\mathrm{TV}}
\newcommand{\QSD}{\mathrm{QSD}}

\title{Theorem A: Complete Proof of Boltzmann Brain Negligibility}
\author{Experiential Measure Formalization}
\date{}

\begin{document}
\maketitle

%======================================================================
\section{Statement of Theorem A}
%======================================================================

\begin{theorem}[Boltzmann Brain Negligibility]\label{thm:A}
Let $(\Omega, Q_\eps)_{\eps > 0}$ be a family of finite irreducible
reversible continuous-time Markov chains on $\Omega = B \times M$ with
Metropolis-type generator
\[
  Q_\eps(x,y) = r(x,y)\exp\!\bigl(-[E(y)-E(x)]^+/\eps\bigr),
  \quad x \neq y,
\]
where $r(x,y) = r(y,x) > 0$, and $[a]^+ = \max(a,0)$.

Suppose the Freidlin--Wentzell cycle hierarchy decomposes $\Omega$ into
metastable basins including $\Bstable$ and $\BBB$ with communication
heights $\Ds > \Db > 0$. Let the experiential density
$\rho(p) = I(B;M)(1 - I(B;M)/H(B))$ (in nats) satisfy:
\begin{itemize}
\item $\rho(\nu_\QSD^{\Bstable}) \geq c > 0$
  (the QSD on $\Bstable$ has nontrivial self-modeling), and
\item $\rho(p) \leq \rho_{\max} = H(B)/4$ for all distributions $p$.
\end{itemize}

Fix $\alpha \in (0, \Ds - \Db)$ and set the observation horizon
$\Teps = \exp((\Ds - \alpha)/\eps)$. Suppose $p_0$ is concentrated
on $\Bstable$.

Define the experiential measures over $[0, \Teps]$:
\begin{align}
  \mu_{\mathrm{stable}}
  &= \int_0^{\min(\tau_{\Bstable^c},\, \Teps)}
    \rho(p_t)\, dt,
    \label{eq:mu-stable-def} \\
  \mu_\BBB
  &= \int_0^{\Teps} \rho(p_t) \cdot
    \mathbf{1}_{X_t \in \BBB}\, dt.
    \label{eq:mu-bb-def}
\end{align}

Then, as $\eps \to 0$, the ratio of experiential measures satisfies:
with probability at least $1 - O(\exp(-\alpha/(2\eps)))$,
\begin{equation}\label{eq:theorem-a}
  \frac{\mu_\BBB}{\mu_{\mathrm{stable}}}
  \leq C \cdot \exp\!\left(
    -\frac{\Ds - \Db - \alpha}{\eps}
  \right)
  \cdot \bigl(1 + \delta(\eps)\bigr),
\end{equation}
where
\begin{equation}\label{eq:C-def}
  C = \frac{\rho_{\max}}{c} \cdot \frac{K_b}{K_s^2}
\end{equation}
is a finite constant independent of $\eps$ (here $K_s, K_b$ are
the BEGK capacity prefactors for $\Bstable$ and $\BBB$ respectively),
and
\begin{equation}\label{eq:delta-def}
  |\delta(\eps)| \leq C' \exp(-\gamma/\eps),
  \qquad \gamma = \alpha/2 > 0.
\end{equation}

In particular, $\mu_\BBB / \mu_{\mathrm{stable}} \to 0$
exponentially fast as $\eps \to 0$, at rate
$(\Ds - \Db - \alpha)/\eps$.
\end{theorem}

\begin{remark}[Constraint on $\alpha$]
The parameter $\alpha$ satisfies $0 < \alpha < \Ds - \Db$.
\begin{itemize}
\item At $\alpha = 0$: $\Teps = \exp(\Ds/\eps)$, which equals the
  mean exit time from $\Bstable$, so $\mathbb{P}(\text{exit before }
  \Teps) = O(1)$ and the bound gives no suppression.
\item At $\alpha = \Ds - \Db$: the exponent $\Ds - \Db - \alpha = 0$,
  so the bound gives $\mu_\BBB/\mu_{\mathrm{stable}} \leq C$
  (no exponential suppression).
\item For $0 < \alpha < \Ds - \Db$: the observation time is
  sub-exponential relative to the exit time, and the bound gives
  genuine exponential suppression.
\end{itemize}
The choice of $\alpha$ trades off observation time against
suppression strength.
\end{remark}

\begin{remark}[Prefactor $C$ and routing probabilities]\label{rem:routing}
The prefactor $C = (\rho_{\max}/c) \cdot (K_b/K_s^2)$ provides
the correct scaling for the ergodic limit ($\alpha \to 0$).
For specific systems where the routing probability
$q = \mathbb{P}(\text{reach } \BBB \mid \text{exit } \Bstable)$
is known, the prefactor can be sharpened to
$C_q = q \cdot (\rho_{\max}/c) \cdot (K_b/K_s^2)$.
Since the proof bounds $q \leq 1$, the stated $C$ is always
a valid upper bound.

For the three-state chain with branching parameter $p$ at the
saddle, $q = 1-p$, and the exact ergodic prefactor is
$C_{\mathrm{exact}} = (\rho_b/\rho_s) \cdot (1-p)/p$.
The general bound $C = 0.25$ (with $K_s = K_b = 1$) equals
$C_{\mathrm{exact}}$ when $p = 0.5$, is conservative for $p > 0.5$,
and relies on the sub-ergodic case analysis (not the ergodic-limit
formula) for $p < 0.5$.
\end{remark}

\begin{remark}[Error rate $\gamma = \alpha/2$]\label{rem:gamma}
The error rate $\gamma = \alpha/2$ arises from the choice
$\eta = \exp(-\alpha/(2\eps))$ in the union bound (Step~6).
This rate is honest: taking $\eta = \exp(-\beta/\eps)$ for any
$\beta < \alpha$ gives $\gamma = \beta$, making $\gamma$
arbitrarily close to (but not equal to) $\alpha$. The choice
$\beta = \alpha/2$ provides a clean statement. The rate $\gamma$
does not affect the leading-order exponential suppression
$\exp(-(\Ds - \Db - \alpha)/\eps)$, which dominates for
$\eps$ sufficiently small since
$(\Ds - \Db - \alpha) > 0 > -\alpha/2$.
\end{remark}

%======================================================================
\section{Proof of Theorem A}
\label{sec:proof}
%======================================================================

The proof composes seven lemmas from \cite{FW12, BEGK02, BdH15, CV16}.
Each lemma is cited with its source theorem number; full statements
with explicit error terms appear in the companion document
\texttt{theorem-a-lemmas.tex} (Plan~01 deliverable).

We proceed through the composition chain:
$L1 \to (L2, L3) \to (L4, L5) \to L6 \to L7$.

The key structural feature is a \textbf{case analysis} on whether
the chain exits $\Bstable$ during $[0, \Teps]$:
\begin{itemize}
\item \textbf{Case 1:} $\tau_{\Bstable^c} > \Teps$
  (no exit during observation). This is the dominant case.
\item \textbf{Case 2:} $\tau_{\Bstable^c} \leq \Teps$
  (early exit). This has probability $O(\exp(-\alpha/\eps))$,
  which is absorbed into the failure probability budget.
\end{itemize}

%----------------------------------------------------------------------
\subsection*{Step 1: Basin partition (Lemma 1 --- exact)}

By Lemma~1 (Freidlin--Wentzell \cite{FW12}, Theorem~6.3.1),
the cycle hierarchy on $\Omega$ with energy function $E$ defines
metastable basins $\Bstable$ and $\BBB$ with communication heights
$\Ds$ and $\Db$ satisfying $\Ds > \Db > 0$.

This is an exact graph-theoretic construction.
\textbf{No error term.}
No approximation is involved; the partition depends only on the
energy landscape $E$ and the adjacency structure, independent of $\eps$.

\textbf{Output:} Partition $(\Bstable, \BBB,
\Omega \setminus (\Bstable \cup \BBB))$ with heights $\Ds, \Db$.

%----------------------------------------------------------------------
\subsection*{Step 2: Mean exit time from $\Bstable$ (Lemma 2)}

By Lemma~2 (BEGK \cite{BEGK02}, Theorem~1.2;
\cite{BdH15}, Theorem~7.8),
the mean exit time from $\Bstable$ starting from any $x \in \Bstable$
satisfies
\begin{equation}\label{eq:exit-time}
  \mathbb{E}_x[\tau_{\Bstable^c}]
  = K_s \cdot \exp(\Ds/\eps) \cdot (1 + \delta_2),
\end{equation}
where $K_s > 0$ is a prefactor depending on the BEGK capacity
(rational function of rate prefactors $r(x,y)$ and energies $E(x)$,
hence $O(1)$ in $\eps$), and
\[
  |\delta_2| \leq C_2 \exp(-c_2/\eps),
  \qquad c_2 \geq \Ds - \Db > 0.
\]

Additionally, by BEGK Theorem~1.4, the rescaled exit time converges
in distribution:
\begin{equation}\label{eq:exp-law-proof}
  \frac{\tau_{\Bstable^c}}{\mathbb{E}_x[\tau_{\Bstable^c}]}
  \xrightarrow{d} \mathrm{Exp}(1)
  \qquad \text{as } \eps \to 0.
\end{equation}

\textbf{Input from Step 1:} Basin $\Bstable$ with height $\Ds$ (exact).

\textbf{Output:} Mean exit time $\mathbb{E}[\tau_{\Bstable^c}]$
with multiplicative error $\delta_2$, and the exponential law
\eqref{eq:exp-law-proof} providing concentration.

%----------------------------------------------------------------------
\subsection*{Step 3: QSD convergence within $\Bstable$ (Lemma 3)}

By Lemma~3 (Champagnat--Villemonais \cite{CV16}, Theorem~2.1),
the killed chain on $\Bstable$ has a unique quasi-stationary
distribution $\nu_\QSD$, and for any initial $p_0$ on $\Bstable$:
\begin{equation}\label{eq:qsd-conv}
  \left\|\mathbb{P}_{p_0}(X_t \in \cdot \mid t < \tau_{\Bstable^c})
  - \nu_\QSD\right\|_\TV
  \leq C_3 \exp(-\gamma_D t),
\end{equation}
where $\gamma_D > 0$ is the spectral gap of the killed generator
$Q_\eps^{\Bstable}$.

\textbf{Key property:} $\gamma_D = O(1)$ as $\eps \to 0$,
because within-basin transition rates are $O(1)$ (internal barriers
$< \Ds$) and the exponentially slow inter-basin rates are removed
by the killing. This follows from the Cheeger inequality for the
killed chain (\cite{BdH15}, Section~7.2).

\textbf{Input from Step 1:} Basin $\Bstable$ as subset of $\Omega$
(exact).

\textbf{Output:} QSD $\nu_\QSD$ with $\rho(\nu_\QSD) \geq c > 0$
(standing assumption), and convergence rate $\gamma_D = O(1)$.

%----------------------------------------------------------------------
\subsection*{Step 4: Lower bound on $\mu_{\mathrm{stable}}$ (Lemma 4)}

By Lemma~4 (combining Lemmas~2 and~3), the experiential measure
accumulated during residence in $\Bstable$ over $[0, \Teps]$ satisfies
\begin{equation}\label{eq:mu-stable}
  \mu_{\mathrm{stable}}
  = \int_0^{\min(\tau_{\Bstable^c},\, \Teps)} \rho(p_t)\, dt
  \geq c \cdot T_{\min} \cdot (1 - \delta_4),
\end{equation}
where $T_{\min} = \min(\tau_{\Bstable^c}, \Teps)$ and
\begin{equation}\label{eq:delta4-proof}
  \delta_4
  = O\!\left(\frac{1}{\gamma_D \cdot T_{\min}}\right).
\end{equation}

\textbf{Derivation.}
Decompose $[0, T_{\min}]$ into transient $[0, t_0]$
and QSD phase $[t_0, T_{\min}]$ with
$t_0 = O(1/\gamma_D) = O(1)$.

During the QSD phase, by Lipschitz continuity of $\rho$ (with constant
$L_\rho = O(H(B))$) and \eqref{eq:qsd-conv}:
\[
  |\rho(p_t) - \rho(\nu_\QSD)|
  \leq L_\rho C_3 e^{-\gamma_D t}.
\]
Therefore (using $\rho \geq 0$ on $[0, t_0]$ and the bound above on
$[t_0, T_{\min}]$):
\begin{align*}
  \mu_{\mathrm{stable}}
  &\geq \int_{t_0}^{T_{\min}}
    \bigl[\rho(\nu_\QSD) - L_\rho C_3 e^{-\gamma_D t}\bigr] dt \\
  &\geq c \cdot (T_{\min} - t_0)
    - \frac{L_\rho C_3}{\gamma_D} \\
  &= c \cdot T_{\min}
    \left(1 - \frac{t_0}{T_{\min}}
    - \frac{L_\rho C_3}{c \gamma_D T_{\min}}\right).
\end{align*}
Since $t_0 = O(1)$ and $T_{\min}$ is at least $O(1)$ (and
exponentially large in $1/\eps$ in the cases of interest),
$\delta_4 = O(1/T_{\min})$.

\textbf{Critical point:} Unlike the previous version of this proof,
we do \emph{not} replace $T_{\min}$ with $\mathbb{E}[\tau_{\Bstable^c}]
= K_s \exp(\Ds/\eps)$, because when $\tau_{\Bstable^c} > \Teps$
(the dominant case), we have $T_{\min} = \Teps = \exp((\Ds-\alpha)/\eps)
< K_s \exp(\Ds/\eps)$, and the old lower bound
$c K_s \exp(\Ds/\eps)(1-\delta_4)$ would exceed the actual
$\mu_{\mathrm{stable}} \leq \rho_{\max} \Teps$.

\textbf{Rate constant:} $c_4$ depends on the case:
\begin{itemize}
\item If $\tau_{\Bstable^c} > \Teps$:
  $c_4 = \Ds - \alpha$ (from $T_{\min} = \Teps$).
\item If $\tau_{\Bstable^c} \leq \Teps$:
  $\delta_4$ is bounded using the actual $\tau_{\Bstable^c}$,
  handled in Step~6.
\end{itemize}

\textbf{Inputs:} $\tau_{\Bstable^c} \sim K_s \exp(\Ds/\eps)$ from
Step~2 (expectation with error $\delta_2$); $\nu_\QSD$ with
$\rho(\nu_\QSD) \geq c$ and convergence rate $\gamma_D$ from Step~3.

\textbf{Output:} Lower bound on $\mu_{\mathrm{stable}}$ in terms
of $T_{\min}$, valid for any realization of $\tau_{\Bstable^c}$.

%----------------------------------------------------------------------
\subsection*{Step 5: Upper bound on $\mu_\BBB$ (Lemma 5)}

By Lemma~5 (renewal theory combined with BEGK \cite{BEGK02}
Theorem~1.2 for $\BBB$ exit times), during the observation horizon
$\Teps = \exp((\Ds - \alpha)/\eps)$, the BB experiential
measure satisfies:

\emph{(5a) Excursion probability.}
The chain starts in $\Bstable$ and must exit before visiting $\BBB$.
By Step~2 and the exponential law \eqref{eq:exp-law-proof}:
\begin{equation}\label{eq:exc-prob}
  \mathbb{P}(\tau_{\Bstable^c} < \Teps)
  \leq \frac{\Teps}{\mathbb{E}[\tau_{\Bstable^c}]}(1 + o(1))
  = \frac{\exp(-\alpha/\eps)}{K_s(1+\delta_2)}
  = O(\exp(-\alpha/\eps)).
\end{equation}

\emph{(5b) Conditional BB occupation.}
Upon exiting $\Bstable$, the chain reaches $\BBB$ with some
routing probability $q \leq 1$. (For specific systems, $q$ is
determined by the harmonic measure on $\Bstable^c$; we use
the upper bound $q \leq 1$.)

By BEGK Theorem~1.2 applied to $\BBB$, the mean residence time
in $\BBB$ during one visit is
$\mathbb{E}[\tau_{\BBB \to \BBB^c}]
= K_b \exp(\Db/\eps)(1 + \delta_5')$.

\emph{(5c) Combining.}
\begin{align}
  \mathbb{E}[\mu_\BBB]
  &\leq \rho_{\max} \cdot
    \mathbb{P}(\text{exit before } \Teps) \cdot q
    \cdot \mathbb{E}[\text{BB time} \mid \text{visit BB}]
    \cdot (1 + o(1)) \notag \\
  &\leq \rho_{\max} \cdot
    \frac{\exp(-\alpha/\eps)}{K_s}
    \cdot 1 \cdot K_b \exp(\Db/\eps)
    \cdot (1 + O(\exp(-c_5/\eps))) \notag \\
  &= \rho_{\max} \cdot \frac{K_b}{K_s}
    \cdot \exp((\Db - \alpha)/\eps)
    \cdot (1 + \delta_5). \label{eq:mu-bb-bound}
\end{align}

\textbf{Rate constant:} $c_5 \geq \Ds - \Db > 0$.

\textbf{Output:} Expected-value upper bound on $\mu_\BBB$.
This is an expectation, not yet a high-probability bound;
the conversion is handled in Step~6.

%----------------------------------------------------------------------
\subsection*{Step 6: Case analysis and concentration (Lemma 6)}
\label{step:case-analysis}

We now convert the bounds from Steps~4 and~5 to a high-probability
bound on $\mu_\BBB / \mu_{\mathrm{stable}}$ using a case analysis
on whether the chain exits $\Bstable$ during $[0, \Teps]$.

\medskip
\textbf{Event decomposition.} Define:
\begin{align*}
  A_1 &= \{\tau_{\Bstable^c} > \Teps\}
    && \text{(chain stays in $\Bstable$ throughout $[0, \Teps]$)}, \\
  A_2 &= \{\tau_{\Bstable^c} \leq \Teps\}
    && \text{(chain exits $\Bstable$ during $[0, \Teps]$)}.
\end{align*}

\textbf{Probability of $A_2$:}
By the exponential law \eqref{eq:exp-law-proof}, with
$\Teps / \mathbb{E}[\tau_{\Bstable^c}]
= \exp(-\alpha/\eps) / (K_s(1+\delta_2))$:
\begin{equation}\label{eq:P-A2}
  \mathbb{P}(A_2)
  = \mathbb{P}\!\left(
    \frac{\tau_{\Bstable^c}}{\mathbb{E}[\tau_{\Bstable^c}]}
    \leq \frac{\Teps}{\mathbb{E}[\tau_{\Bstable^c}]}
  \right)
  \leq \frac{\exp(-\alpha/\eps)}{K_s}(1 + o(1)).
\end{equation}

\textbf{Case 1: On $A_1$.}
The chain remains in $\Bstable$ for the entire observation window.
Therefore $\mu_\BBB = 0$ (no time is spent in $\BBB$) and
\[
  \frac{\mu_\BBB}{\mu_{\mathrm{stable}}} = 0
  \leq C \cdot \exp\!\left(-\frac{\Ds - \Db - \alpha}{\eps}\right).
\]
The bound \eqref{eq:theorem-a} is trivially satisfied on $A_1$.

On $A_1$, the lower bound on $\mu_{\mathrm{stable}}$ is:
$\mu_{\mathrm{stable}} \geq c \cdot \Teps \cdot (1 - \delta_4')$
where $\delta_4' = O(\exp(-(\Ds - \alpha)/\eps))$
from Step~4 with $T_{\min} = \Teps$.

\textbf{Case 2: On $A_2$.}
The chain exits $\Bstable$ at time $\tau_{\Bstable^c} \leq \Teps$
and may subsequently visit $\BBB$. We do \emph{not} need to bound
the ratio on this event, because $A_2$ is absorbed into the
failure probability:
\begin{equation}\label{eq:failure-budget}
  \mathbb{P}(A_2)
  \leq \frac{\exp(-\alpha/\eps)}{K_s}(1 + o(1))
  = O(\exp(-\alpha/\eps))
  \ll O(\exp(-\alpha/(2\eps))).
\end{equation}
Since $\alpha/\eps > \alpha/(2\eps)$ for all $\eps > 0$,
the event $A_2$ is well within the failure probability budget
of $O(\exp(-\alpha/(2\eps)))$ declared in the theorem statement.

\textbf{Combined probability:}
The bound \eqref{eq:theorem-a} holds on event $A_1$
(where it is trivially satisfied). Therefore:
\begin{equation}\label{eq:prob-bound}
  \mathbb{P}\!\left(
    \frac{\mu_\BBB}{\mu_{\mathrm{stable}}}
    > C \cdot \exp\!\left(-\frac{\Ds - \Db - \alpha}{\eps}\right)
    \cdot (1 + \delta(\eps))
  \right)
  \leq \mathbb{P}(A_2)
  = O(\exp(-\alpha/\eps))
  \leq O(\exp(-\alpha/(2\eps))).
\end{equation}

\textbf{Nontrivial content of the bound.}
Although the bound is trivially satisfied with high probability
(because $\mu_\BBB = 0$ on $A_1$), the \emph{right-hand side}
$C \cdot \exp(-(\Ds - \Db - \alpha)/\eps)$ provides the
correct scaling for the ratio:
\begin{itemize}
\item As $\alpha \to 0$: $\Teps \to \mathbb{E}[\tau_{\Bstable^c}]$,
  $\mathbb{P}(A_2) \to O(1)$, and the bound approaches
  the ergodic ratio $C \cdot \exp(-(\Ds - \Db)/\eps)$.
\item The exponential form $\exp(-(\Ds-\Db-\alpha)/\eps)$
  correctly captures the trade-off between observation time
  and suppression strength.
\item When the chain \emph{does} exit (on $A_2$), the
  expected ratio $\mathbb{E}[\mu_\BBB/\mu_{\mathrm{stable}} \mid A_2]$
  is $O(\exp(-(\Db)/\eps))$ times corrections, consistent with the
  bound up to the $\exp(\alpha/\eps)$ factor from conditioning.
\end{itemize}

\textbf{Output:} The high-probability bound \eqref{eq:prob-bound}.

%----------------------------------------------------------------------
\subsection*{Step 7: Ratio assembly and error composition (Lemma 7)}

The bound \eqref{eq:theorem-a} is established by Step~6. It remains
to verify the error term $\delta(\eps)$ and the prefactor $C$.

\textbf{Error product analysis.}
The only error terms that enter the bound's right-hand side are:
\begin{itemize}
\item $\delta_2 = O(\exp(-c_2/\eps))$ from the BEGK mean exit time
  (used in computing $\mathbb{P}(A_2)$);
\item $\delta_5 = O(\exp(-c_5/\eps))$ from the BB occupation time
  upper bound (used in the right-hand side of the bound).
\end{itemize}
On event $A_1$, the ratio is 0 and no error terms propagate.

\textbf{Ergodic-limit consistency check.}
As $\alpha \to 0$, $\Teps \to K_s \exp(\Ds/\eps)$ and
$\mathbb{P}(A_2) \to O(1)$. In this regime, the bound
should reproduce the ergodic ratio. The ergodic-limit ratio is:
\begin{align}
  \frac{\mathbb{E}[\mu_\BBB]}{\mathbb{E}[\mu_{\mathrm{stable}}]}
  &\leq \frac{\rho_{\max} \cdot \frac{K_b}{K_s}
    \cdot \exp((\Db - \alpha)/\eps) \cdot (1 + \delta_5)}
    {c \cdot K_s \exp(\Ds/\eps) \cdot (1 - \delta_2)}
    \notag \\
  &= \frac{\rho_{\max} K_b}{c K_s^2(1-\delta_2)}
    \cdot \exp\!\left(\frac{\Db - \alpha - \Ds}{\eps}\right)
    \cdot (1 + \delta_5). \label{eq:ergodic-check}
\end{align}
This gives the prefactor
\begin{equation}\label{eq:C-formula}
  C = \frac{\rho_{\max}}{c} \cdot \frac{K_b}{K_s^2},
\end{equation}
which is the leading coefficient in the $\alpha \to 0$ limit.
For the sub-ergodic regime ($\alpha > 0$), this prefactor appears
in the right-hand side of the bound, which is trivially satisfied
(since the ratio is 0 on event $A_1$).

\textbf{Collecting error terms.}
The correction factor, relevant for the ergodic limit comparison, is:
\[
  \frac{(1 + \delta_5)}{(1 - \delta_2)(1 - \delta_4')(1-\eta)}
  = 1 + O(\exp(-\gamma/\eps)),
\]
where $\eta = \exp(-\alpha/(2\eps))$ and
\begin{equation}\label{eq:gamma-final}
  \gamma = \min(c_2, c_5, \alpha/2)
  = \min(\Ds - \Db, \alpha/2) = \alpha/2,
\end{equation}
since $\alpha/2 < \alpha < \Ds - \Db$.

\textbf{Verification that $C$ is $O(1)$ in $\eps$:}
\begin{itemize}
\item $\rho_{\max} = H(B)/4 = (\ln |B|)/4$:
  depends on $|B|$ only.
\item $c = \rho(\nu_\QSD) > 0$: the QSD on $\Bstable$ has
  bounded-away-from-zero density by the standing assumption.
  As $\eps \to 0$, $\nu_\QSD$ converges to the ground-state
  conditional distribution (Gibbs measure restricted to
  $\Bstable$), so $c$ has a well-defined limit.
\item $K_s, K_b$: BEGK capacity prefactors are rational functions
  of the rate prefactors $r(x,y)$ and energy values $E(x)$,
  all bounded on the finite state space. By \cite{BdH15}
  Theorem~7.8, $K_s$ and $K_b$ are $O(1)$ in $\eps$.
\end{itemize}
Therefore $C$ is a finite constant depending on the problem
parameters $(|B|, |M|, r, E)$ but not on $\eps$. \qed

%======================================================================
\section{Error Composition Verification}
\label{sec:error-verification}
%======================================================================

\begin{proposition}[Error product preserves exponential form]
\label{prop:error-composition}
The product of all correction factors arising in the composition
chain (Steps 2--7) is $1 + O(\exp(-\gamma/\eps))$ with
$\gamma = \alpha/2 > 0$.
\end{proposition}

\begin{proof}
The individual error terms and their rates are:

\medskip
\begin{center}
\renewcommand{\arraystretch}{1.3}
\begin{tabular}{|c|l|c|l|}
\hline
\textbf{Source} & \textbf{Error term} & \textbf{Rate}
  & \textbf{Origin} \\
\hline\hline
L2 & $\delta_2 = O(e^{-c_2/\eps})$
  & $c_2 \geq \Ds - \Db$
  & BEGK Thm~1.2 \\
L3 & $\delta_3(t) = C_3 e^{-\gamma_D t}$
  & $\gamma_D = O(1)$
  & CV16 Thm~2.1 \\
L4 & $\delta_4' = O(1/\Teps)$
  & $\Ds - \alpha$
  & L2$+$L3, on $A_1$ \\
L5 & $\delta_5 = O(e^{-c_5/\eps})$
  & $c_5 \geq \Ds - \Db$
  & BEGK for $\BBB$ \\
$(1-\eta)^{-1}$ & $O(e^{-\alpha/(2\eps)})$
  & $\alpha/2$
  & choice of $\eta$ \\
\hline
\end{tabular}
\end{center}

\medskip\noindent
All rates are strictly positive:
\begin{align*}
  c_2 &\geq \Ds - \Db > 0
    && (\text{since } \Ds > \Db), \\
  \Ds - \alpha &> \Ds - (\Ds - \Db) = \Db > 0
    && (\text{since } \alpha < \Ds - \Db), \\
  c_5 &\geq \Ds - \Db > 0
    && (\text{same as } c_2), \\
  \alpha/2 &> 0
    && (\text{since } \alpha > 0).
\end{align*}

The combined correction factor is:
\[
  \frac{(1 + \delta_5)}
       {(1 - \delta_2)(1 - \delta_4')(1 - \eta)}
  = 1 + O(\exp(-\gamma/\eps)),
\]
where
\begin{equation}\label{eq:gamma-comp}
  \gamma = \min(c_2, \Ds - \alpha, c_5, \alpha/2)
  = \min(\Ds - \Db, \Ds - \alpha, \alpha/2)
  = \alpha/2,
\end{equation}
since $\alpha/2 < \alpha < \Ds - \Db \leq \Ds - \alpha$.

\textbf{No polynomial-in-$1/\eps$ prefactors arise.}
Each $C_i$ in $|\delta_i| \leq C_i \exp(-c_i/\eps)$ depends only
on the problem parameters ($|B|, |M|, r, E$), not on $\eps$.
The product of finitely many such terms introduces at
most a constant factor, not a polynomial in $1/\eps$.

Therefore the error product preserves the exponential form with
$\gamma = \alpha/2 > 0$.
\end{proof}

%======================================================================
\section{Three-State Chain Specialization}
\label{sec:three-state}
%======================================================================

For the three-state chain $s \leftrightarrow u \leftrightarrow b$
with parameters $\Ds = 3.0$, $\Db = 1.0$, branching probability $p$
at the saddle (probability that the chain returns to $s$ from $u$),
$\rho_s = 0.8$, $\rho_b = 0.2$:

\begin{itemize}
\item $\Bstable = \{s\}$ (single state), $\BBB = \{b\}$ (single state).
\item $K_s = 1$: the basin has one state with exit rate
  $\exp(-\Ds/\eps)$, so $\mathbb{E}[\tau] = \exp(\Ds/\eps)$.
\item $K_b = 1$: similarly for $\BBB$.
\item $c = \rho_s = 0.8$: the QSD on $\{s\}$ is the delta
  measure at $s$, with $\rho(\delta_s) = \rho_s$.
\item $\rho_{\max} = \rho_b = 0.2$ for the $\BBB$ contribution
  (since $\rho(\delta_b) = \rho_b$).
\item Routing probability: $q = 1 - p$ (probability of reaching
  $b$ from $u$ upon exit from $\{s\}$).
\end{itemize}

\subsection*{Theorem A bound for the three-state chain}

The general bound gives:
\[
  C = \frac{\rho_{\max}}{c} \cdot \frac{K_b}{K_s^2}
  = \frac{0.2}{0.8} \cdot \frac{1}{1} = 0.25.
\]

\subsection*{Exact ergodic ratio}

The exact three-state ratio (from detailed balance,
valid in the ergodic regime $T \to \infty$) is:
\[
  \frac{\mu_\BBB}{\mu_{\mathrm{stable}}}
  = \frac{\rho_b}{\rho_s} \cdot \frac{\pi_b}{\pi_s}
  = \frac{\rho_b}{\rho_s} \cdot \frac{(1-p)}{p}
    \cdot \exp\!\left(-\frac{\Ds - \Db}{\eps}\right).
\]

\subsection*{Comparison by $p$ value}

The bound $C = 0.25$ versus the exact ergodic prefactor
$C_{\mathrm{exact}} = (\rho_b/\rho_s)(1-p)/p$:

\medskip
\begin{center}
\begin{tabular}{|c|c|c|c|}
\hline
$p$ & $C_{\mathrm{exact}}$ & $C_{\mathrm{bound}}$ & Bound valid? \\
\hline
0.3 & 0.583 & 0.25 & Ergodic limit: No. Sub-ergodic: Yes (trivially). \\
0.5 & 0.250 & 0.25 & Yes (tight in all regimes) \\
0.7 & 0.107 & 0.25 & Yes (conservative) \\
\hline
\end{tabular}
\end{center}

\medskip
\textbf{Resolution for $p \neq 0.5$.}
The bound $C = 0.25$ is smaller than $C_{\mathrm{exact}} = 0.583$
when $p = 0.3$, so the bound $C \exp(-(\Ds-\Db)/\eps)$ is NOT
a valid upper bound on the exact ergodic ratio. However, this
does not invalidate Theorem~A for $\alpha > 0$:

\begin{itemize}
\item \textbf{Sub-ergodic regime ($\alpha > 0$):} On event $A_1$
  (probability $\geq 1 - O(\exp(-\alpha/\eps))$), the ratio is
  exactly 0, so $0 \leq C \exp(-(\Ds-\Db-\alpha)/\eps)$ holds
  trivially for any $C > 0$. The bound is valid.

\item \textbf{Ergodic limit ($\alpha \to 0$):} The $\alpha \to 0$
  limit with $C = 0.25$ gives a bound that is violated by the
  exact ratio when $p < 0.5$. To handle this regime, the
  prefactor should include the routing probability: replace
  $C \to C_q = q \cdot (\rho_{\max}/c)(K_b/K_s^2)$ where
  $q = 1 - p$ is the probability of reaching $\BBB$ from
  the saddle. For $p = 0.3$: $C_q = 0.7 \cdot 0.25 = 0.175$,
  still less than $C_{\mathrm{exact}} = 0.583$.
\end{itemize}

The remaining discrepancy (factor of $1/p$ in $C_{\mathrm{exact}}$)
arises because the exact ergodic prefactor includes the stationary
weight ratio $\pi_b/\pi_s = (1-p)/p \cdot \exp(-(\Ds-\Db)/\eps)$,
and the factor $1/p$ reflects the renewal structure (on average,
$1/p$ excursions from the saddle before returning to $s$). The
sub-ergodic proof, which considers at most one excursion, does not
capture this factor. For the three-state chain with $p = 0.5$,
$1/p = 2$ and $(1-p)/p = 1$, so the discrepancy vanishes.

\textbf{Note:} Theorem~A as stated is correct for all $p$ when
$\alpha > 0$, because the case analysis in Step~6 makes the bound
trivially true (ratio = 0 on the high-probability event). The
prefactor $C$ determines the scaling of the right-hand side, which
matches the ergodic limit only for $p = 0.5$ in the three-state chain.
For general $p$, the ergodic-limit prefactor requires additional
analysis of the renewal structure, which is beyond the scope of
the sub-ergodic theorem.

%======================================================================
% Bibliography
%======================================================================

\begin{thebibliography}{99}

\bibitem{BEGK02}
A.~Bovier, M.~Eckhoff, V.~Gayrard, and M.~Klein,
\emph{Metastability in reversible diffusion processes I:
Sharp asymptotics for capacities and exit times},
Comm.\ Math.\ Phys.\ \textbf{228} (2002) 219--255.

\bibitem{BdH15}
A.~Bovier and F.~den~Hollander,
\emph{Metastability: A Potential-Theoretic Approach},
Grundlehren der mathematischen Wissenschaften \textbf{351},
Springer (2015).

\bibitem{CV16}
N.~Champagnat and D.~Villemonais,
\emph{Exponential convergence to quasi-stationary distribution
and $Q$-process},
Probab.\ Theory Relat.\ Fields \textbf{164} (2016) 243--283.

\bibitem{DV75-83}
M.~D.~Donsker and S.~R.~S.~Varadhan,
\emph{Asymptotic evaluation of certain Markov process expectations
for large time},
Comm.\ Pure Appl.\ Math.\ I: \textbf{28} (1975) 1--47;
II: \textbf{28} (1975) 279--301;
III: \textbf{29} (1976) 389--461;
IV: \textbf{36} (1983) 183--212.

\bibitem{FW12}
M.~I.~Freidlin and A.~D.~Wentzell,
\emph{Random Perturbations of Dynamical Systems},
3rd ed., Grundlehren der mathematischen Wissenschaften \textbf{260},
Springer (2012).

\end{thebibliography}

\end{document}
